\newpage
\thispagestyle{empty}

%

% Define block styles
\definecolor{StartEnd}{RGB}{251,232,166}
\definecolor{Question}{RGB}{48,60,108}
\definecolor{Action}{RGB}{244,151,108}
\definecolor{Border}{RGB}{209,253,254}

\tikzstyle{thickBlock} = [rectangle, draw, fill={StartEnd}, 
    text width=12em, text centered, ultra thick, minimum height=2em]
\tikzstyle{block} = [rectangle, draw, fill={Action}, 
    text width=12em, text centered, thick, minimum height=2em]
\tikzstyle{line} = [draw, -stealth', line width=0.05cm]
\tikzstyle{cloud} = [draw={Question}, rectangle, thick, rounded corners=8, fill={Question}, minimum height=2em, text width=12em, text badly centered,text=Border ]

\mbox{}
\vfill
\begin{minipage}[c]{.45\textwidth}
\centering
  \includegraphics[height=80mm,center]{img/Frits.png}
  \captionof{figure}{\Large{\texttt{Frits}}}
\end{minipage}%
\hfill
\begin{minipage}[c]{.45\textwidth}
  \centering
  \includegraphics[height=80mm,center]{img/Willem.png}
  \captionof{figure}{\Large{\texttt{Willem}}}
\end{minipage}
\vfill
\begin{figure}[H]
   \centering
    \begin{tikzpicture}[node distance=1cm and 1cm, >=stealth, bend angle=45,auto]
    % Place nodes
    \node [thickBlock] (block_1) {\large{\textbf{Start beurt}}};
    \node [cloud, right=of block_1] (cloud_1) {Heb je \'e\'en of meer kaarten?};
    \node [thickBlock, right=of cloud_1] (block_end_1) {\large{\textbf{Einde beurt}}};
    \node [cloud, below=of cloud_1] (cloud_2) {Kan \'en wil je een zet\footnotemark[1] doen?};
    \node [block, below=of cloud_2] (block_2) {Doe een zet\footnotemark[1]};
    \node [cloud, right=of block_2] (cloud_3) {Heb je \'e\'en kaart?};
    \node [cloud, below=of cloud_3] (cloud_4) {Heb je geen kaarten?};
    \node [block, above=of cloud_3] (block_3) {Zeg \quotes{laatste Frits}};
    \node [block, left=of cloud_4] (block_4) {Zeg \quotes{Frits} en \textbf{Frits}\footnotemark[2]};
    \node [block, left=of cloud_2] (block_5) {\textbf{Frits}\footnotemark[2] en pak twee kaarten\footnotemark[3]};
    \node [cloud, below=of block_5] (cloud_5) {Kan \'en wil je een zet\footnotemark[1] doen?};
    \node [block, below=of cloud_5] (block_6) {Begin een nieuwe stapel};
    \node [thickBlock, below=of block_4] (block_end_2) {\large{\textbf{Einde beurt}}};
    % Draw edges
     \path [line] (block_1) -/ (cloud_1);
     \path [line] (cloud_1) -/ node {\textbf{Ja}} (cloud_2);
     \path [line] (cloud_1) -- node[above] {\textbf{Nee}} (block_end_1);
     \path [line] (cloud_2) -/ node {\textbf{Ja}} (block_2);
     \path [line] (block_2) -/ (cloud_3);
     \path [line] (cloud_3) -/ node[right] {\textbf{Ja}} (block_3);
     \path [line] (cloud_3) -/ node {\textbf{Nee}} (cloud_4);
     \path [line] (cloud_4) -/ node[above] {\textbf{Ja}} (block_4);
     \path [line] (block_3) -/ (block_end_1);
     \path [line] (cloud_2) -/ node[above] {\textbf{Nee}} (block_5);
     \path [line] (block_5) -/ (cloud_5);
     \path [line] (cloud_5) -/ node[above] {\textbf{Ja}} (block_2);
     \path [line] (cloud_5) -/ node {\textbf{Nee}} (block_6);
     \path [line] (block_6) |- (block_end_2);
     \path [line] (block_4) -- (block_end_2);
     \path [line] (cloud_4) |- node[right] {\textbf{Nee}} (block_end_2);
     \node[align=center,font=\bfseries, yshift=3em] (title) at (current bounding box.north){\Huge{\texttt{Een Fritsbeurt}}};
\end{tikzpicture}
\end{figure}
\vfill


\footnotetext[1]{zie pagina \pageref{sec:regels_kort} en \pageref{zetkort:caroline} voor de zetten}
\footnotetext[2]{\textbf{Fritsen} is het drinken van één shotglaasje bier (zie definitie \ref{def:fritsen} op pagina \pageref{def:fritsen})}
\footnotetext[3]{zolang er dichte kaarten zijn (zie regels \ref{item:geen_kaart_1} en \ref{item:geen_kaart_2} op pagina \pageref{item:geen_kaart_1})}
\newpage

